\clearpage
\renewcommand{\sectioncolor}{indigo}
\section{The bridge to Part IX}
\markboth{The bridge}{}

Everything in §5 and §6 generalises, and the generalisation is exactly what Part IX was about.

\begin{center}\small
\begin{tabular}{@{}>{\raggedright\arraybackslash}p{7.6cm}>{\raggedright\arraybackslash}p{8.9cm}@{}}
\toprule
\rowcolor{indigo!12}
\textbf{\color{indigo}Here --- 1D Ising} & \textbf{\color{indigo}There --- molecular systems (Part IX)}\\
\midrule
$T$, a $2\times2$ \textbf{transfer matrix} & $\mathcal{P}_t$, the \textbf{transfer operator} on conformation space\\
$Z=\operatorname{Tr}(T^N)$, a sum over $2^N$ states & $Z=\int e^{-\beta U(x)}dx$, uncountably many states\\
largest eigenvalue $\lambda_+$ $\Rightarrow$ free energy & leading eigenvector $\Rightarrow$ the \textbf{equilibrium measure} $\mu$\\
gap $\lambda_+/\lambda_-$ $\Rightarrow$ correlation length $\xi$ & remaining eigenvalues $\Rightarrow$ \textbf{relaxation timescales} $t_i=-\tau/\ln\lambda_i$\\
eigenvectors of a $2\times2$ matrix & eigenvectors $\Rightarrow$ the \textbf{slow collective variables} (what VAMPnets learn)\\
\rowcolor{alert!7}
MCMC fails when $\xi\gtrsim N$ & MCMC fails when barriers $\gg k_BT$ --- \textbf{the same mechanism}\\
\bottomrule
\end{tabular}
\end{center}

\begin{keyidea}[title={The sentence to keep}]
\begin{center}\itshape\large
It is called a transfer \emph{operator} because it is the continuum limit\\of the transfer
\emph{matrix} you just diagonalised.
\end{center}
\vspace{2pt}
When Part IX said \emph{``the leading eigenvector is the equilibrium measure, and the remaining
eigenvalues give the relaxation timescales''}, that was not an analogy. It is literally Table 2 of
this dossier, with $2\times2$ replaced by infinite-dimensional.
\end{keyidea}

\subsection{What the Ising route does \emph{not} teach you}

Being straight, so you do not over-learn from one example:

\begin{center}\small
\begin{tabular}{@{}>{\raggedright\arraybackslash}p{7.6cm}>{\raggedright\arraybackslash}p{8.9cm}@{}}
\toprule
\rowcolor{alert!10}
\textbf{\color{alert}What it does not show} & \textbf{\color{alert}Where to get it}\\
\midrule
\textbf{Why computing $Z$ is \#P-hard.} In 1D it is not hard --- the transfer matrix is exact. &
 2D (Onsager 1944) and higher. Kardar has the transfer-matrix route into 2D.\\
Phase transitions. 1D short-range systems cannot order at $T>0$. & Chandler ch.\ 5; Kardar on the
 Ising model in $d\ge2$ and the renormalisation group.\\
Continuous configuration space. Here the state space is finite. & Chandler on classical fluids;
 then Part IX §4 on configurational entropy.\\
\bottomrule
\end{tabular}
\end{center}

\clearpage
\section{A twelve-week plan}
\markboth{Twelve-week plan}{}

\begin{center}\small
\begin{tabular}{@{}>{\raggedright\arraybackslash}p{1.5cm}>{\raggedright\arraybackslash}p{5.2cm}>{\raggedright\arraybackslash}p{4.6cm}>{\raggedright\arraybackslash}p{5.0cm}@{}}
\toprule
\rowcolor{bio!14}
\textbf{\color{bio}Weeks} & \textbf{\color{bio}Read} & \textbf{\color{bio}Do} & \textbf{\color{bio}Question to answer in writing}\\
\midrule
\textbf{1--2} & Axler §5, §6, §7.B, §7.C, §10.A &
 Prove $\operatorname{Tr}(T^N)=\lambda_+^N+\lambda_-^N$ from basis-independence of the trace &
 Why does $e^{-\beta\hat H}$ need the spectral theorem to even be defined?\\
\textbf{3} & --- &
 \textbf{Run the companion notebook.} Do exercises 1--3 &
 At what $N$ does brute force become impossible on your laptop? Extrapolate to $N=60$.\\
\textbf{4--5} & Chandler ch.\ 1--3 \newline{\scriptsize (the one book to buy)} &
 Re-derive $f$, $u$, $C$ for the Ising chain without looking &
 Which of the six quantities in Part IX Fig.~2 can you now compute yourself?\\
\textbf{6--7} & Radozycki III: contour integrals, residues &
 Evaluate a Gaussian integral by contour deformation; then by saddle point &
 What \emph{is} the saddle-point approximation approximating?\\
\textbf{8--9} & Kuttler ch.\ 20--21 &
 Write down $\mu(dx)=e^{-\beta U}dx/Z$ as a measure and verify it is a probability measure &
 Where exactly does positivity of $e^{-\beta\hat H}$ get used?\\
\textbf{10} & Kuttler ch.\ 22 (Riesz) &
 --- &
 In what sense does a positive operator \emph{determine} a measure?\\
\textbf{11--12} & Levin \& Peres ch.\ 1--4 \newline{\scriptsize (the second book to buy)} &
 Notebook exercise 4: measure autocorrelation time vs $\xi$ &
 \textbf{Restate §6 of this dossier in the language of mixing times.}\\
\bottomrule
\end{tabular}
\end{center}
\vspace{-2pt}
\begin{center}\begin{minipage}{0.93\textwidth}\footnotesize\color{slate}
\textbf{Table 4.}\; Two books purchased, ten weeks of reading you already own. At the end of week 12
you can read Part IX §7 (Boltzmann generators) and know precisely what problem it is solving.
\end{minipage}\end{center}

\begin{haveit}[title={The companion notebook}]
\bk{Partition\_Function\_Study/Partition\_Function\_01\_Ising\_Transfer\_Matrix.ipynb}\\[5pt]
17 cells --- 9 exposition, 8 code. Runs end to end on \texttt{numpy} alone; plots are optional and it
degrades gracefully if \texttt{matplotlib} is absent. Every cell was executed before delivery.

\begin{center}\small
\begin{tabular}{@{}cl@{}}
\toprule
\rowcolor{bio!12}
\textbf{§} & \textbf{Contents}\\
\midrule
1--2 & the model; deriving the transfer matrix from the Boltzmann weight\\
3 & \textbf{Method B} --- brute force over all $2^N$ states, and the agreement table\\
4 & thermodynamics: $f$, $u$, $C$ from $\ln\lambda_+$\\
5 & the eigenvalue gap and the correlation length\\
6--7 & \textbf{Method C} --- Metropolis, and the ergodicity failure test\\
8 & the bridge table to Part IX\\
9 & six exercises, including ``find the $2^N$ wall'' and the 2D transfer matrix\\
\bottomrule
\end{tabular}
\end{center}
\end{haveit}
